% Source: https://artofproblemsolving.com/community/c3h1565322_circletriangle_geometry_problem
% Stable identity: aops-topic-1565322-q01
Let $O$ be the centre. The triangle fails to contain $O$ in its interior exactly when the three points lie on a closed semicircle (if they lie on a semicircle, the triangle is on one side of its diameter; if not, every half-plane bounded by a line through $O$ contains a vertex, so $O$ is inside). The case with $O$ on a side has probability $0$.

For each point $P_i$, let $E_i$ be the event that the other two points lie on the semicircle starting at $P_i$ and going counterclockwise. Each of the other two points lands there independently with probability $\dfrac{1}{2}$, so $P(E_i)=\dfrac{1}{4}$.

\begin{asy}
usepackage("amsmath");
size(5cm);
defaultpen(fontsize(10pt));
pair O=(0,0);
draw(circle(O,1),linewidth(0.8));
draw(arc(O,1,-20,160),red+linewidth(2));
draw(dir(-20)--dir(160),red+dashed);
dot(O);label("$O$",O,SW);
pair A=dir(-20),B=dir(60),C=dir(130);
draw(A--B--C--cycle,blue+linewidth(0.9));
dot(A);dot(B);dot(C);label("$A$",A,SE);label("$B$",B,NE);label("$C$",C,NW);
label("$B$, $C$ on the semicircle counterclockwise from $A$",(0,-1.3));
label("\mbox{}",(min(currentpicture,user=true)+max(currentpicture,user=true))/2,black); // keeps the text after the figure black when LaTeX typesets the labels inline
\end{asy}

If the three points lie on some semicircle, then the first of them (going counterclockwise along that semicircle) has the other two within half a turn counterclockwise, so some $E_i$ occurs. Two of the events cannot happen together except with probability $0$: if $P_2$ is less than half a turn counterclockwise from $P_1$, then $P_1$ is more than half a turn counterclockwise from $P_2$. Hence
\[
P(\text{the triangle does not contain }O)=P(E_1)+P(E_2)+P(E_3)=\boxed{\dfrac{3}{4}}.
\]
